\subsection{Controlled Perturbation and Pathway Tests}\label{sec:w2l_perturbation}
The experiment drivers apply neural stimuli, external forces, pathway deletions, and body-mass changes to a fixed controller configuration. Replaying all stored controller fields gives 17 left and 16 right strides in the control condition. The corresponding counts are 16/16 after crossed-heel deletion, 13/13 after Ia-chain deletion, 16/14 after Renshaw-cell deletion, 17/15 after heel-cascade deletion, and 13/12 after commissural deletion. Removing Ib load feedback stops stepping in both legs. Thus, Ib feedback is necessary for stepping at this tested operating point; the other deletions retain stepping with varying counts.

The neural stimuli and external-force tests also retain bilateral stepping. Counts alone do not establish balance recovery: the stronger stimulus and both push conditions exceed the \ang{45} tilt criterion. Across ten body-mass conditions, stride counts span 32--38 per leg and the lowest recorded pelvis height is \qty{0.8981}{\meter}. Nine conditions exceed the tilt criterion; one meets the full survival gate. These comparisons demonstrate controlled changes in neural and mechanical conditions while distinguishing continued stepping from postural stability.

\begin{figure}[p]
\centering
\includegraphics[width=\textwidth]{w2l_perturbation_fully_applied.pdf}
\caption[Controlled W2L perturbation and pathway tests]{Supported W2L tests with all stored controller fields applied. The control produces 17/16 left/right strides. Ib deletion stops stepping; commissural deletion retains 13/12 strides. Stimulus and push trials retain stepping, although several exceed the tilt criterion. The controller replay does not constitute a new optimization with the corrected parameter forwarding.}
\label{fig:w2l_perturbation}
\end{figure}
